\documentclass[11pt]{article}

\usepackage{amsmath,amssymb}
\usepackage[margin=1in]{geometry}
\usepackage{booktabs}
\usepackage{array}
\usepackage{hyperref}

\title{Coefficient Stabilization and Invariant-Tracking Feasibility of a Frozen Branch-Local Cochain-Laplacian Hierarchy}
\author{Tomislav Rupi\'c \\ Independent Researcher}
\date{March 2026}

\begin{document}

\maketitle

\begin{abstract}
Phase IX of the frozen HAOS-IIP program asks whether the coefficient-like structures extracted from the already frozen Phase VIII short-time trace regime exhibit reproducible stabilization across refinement and whether simple normalized combinations behave invariant-like on the branch. The phase does not alter Phase V, Phase VI, Phase VII, or the frozen Phase VIII short-time window and probe grid. It reuses the exact deterministic trace proxy on the hierarchy $h = 1 / n_{\mathrm{side}}$ with $n_{\mathrm{side}} \in \{12,24,36,48\}$ and tracks the Phase VIII size-factored quadratic surrogate coefficients $b_0(h)$, $b_1(h)$, and $b_2(h)$. Across refinement, the relative spans are $8.50 \times 10^{-5}$, $0.01522$, and $0.02726$, with shrinking successive coefficient-vector distances $0.647775$, $0.121305$, and $0.042557$. The ratio descriptors $b_1/b_0$ and $b_2/b_0$ are materially more stable than the raw trace coefficients, with relative spans $0.01531$ and $0.02734$. Rescaled descriptors $I_0$, $I_1$, and $I_2$, formed using the dominant Phase VII low-mode exponent, remain inside compact refinement bands with largest span $0.06543$. A deterministic open-grid scalar control hierarchy with the same node counts behaves differently: its ratio spans are larger by a factor $4.44$, its final coefficient-vector distance is larger by a factor $8.11$, and its best simple extrapolation family is linear in $h$ rather than linear in $h^2$. The result is bounded: Phase IX establishes coefficient stabilization and invariant-tracking feasibility for the frozen operator hierarchy. It does not establish continuum coefficients, geometric invariants, universality, or physical correspondence.
\end{abstract}

\section{Scope and Authority}

Phase IX is a stabilization test, not an interpretive extension. The phase consumes four already frozen authority layers:
\begin{itemize}
\item the Phase V recovery-statistics authority;
\item the Phase VI branch-local operator freeze;
\item the Phase VII spectral-feasibility baseline;
\item the Phase VIII spectral-trace contract, including the frozen short-time window and deterministic probe grid.
\end{itemize}

Phase IX is not allowed to modify any of those layers. Its only purpose is to determine whether the coefficient-like objects already exposed by Phase VIII show stabilization, whether normalized combinations behave more regularly than the raw coefficients, and whether the branch can be distinguished from a deterministic generic-graph control.

\section{Frozen Inputs}

The frozen operator remains the branch-local periodic DK2D block cochain Laplacian
\[
\Delta_h,
\]
on the deterministic refinement hierarchy
\[
h = \frac{1}{n_{\mathrm{side}}}, \qquad n_{\mathrm{side}} \in \{12,24,36,48\}.
\]

The frozen trace contract inherited from Phase VIII is:
\begin{itemize}
\item operational short-time window: $\{0.02, 0.03, 0.05, 0.075, 0.1\}$;
\item exact deterministic trace proxy:
\[
T_h(t) = \mathrm{Tr}\!\left(e^{-t \Delta_h}\right);
\]
\item deterministic truncation checkpoints: full spectrum, $\lambda \le 7.5$, and $\lambda \le 7.8$.
\end{itemize}

Phase IX also reuses the dominant Phase VII low-mode exponent
\[
p = 1.972182228379
\]
as the fixed rescaling exponent for its invariant-like descriptors.

The authority references used in this paper are:
\begin{itemize}
\item \url{phase6-operator/phase6_operator_manifest.json};
\item \url{phase7-spectral/phase7_spectral_manifest.json};
\item \url{phase8-trace/phase8_trace_manifest.json};
\item \url{phase9-invariants/phase9_manifest.json};
\item \url{PHASE_IX_COEFFICIENT_CONTRACT_NOTE.md}.
\end{itemize}

\section{Phase IX Probe Set}

Phase IX performs five deterministic checks:
\begin{enumerate}
\item coefficient convergence tracking for the scaled Phase VIII coefficients $b_0(h)$, $b_1(h)$, and $b_2(h)$;
\item ratio stabilization for
\[
\frac{b_1(h)}{b_0(h)}, \qquad \frac{b_2(h)}{b_0(h)};
\]
\item rescaled descriptor tracking for
\[
I_k(h) = h^p a_k(h),
\]
where $a_k(h)$ are the raw quadratic coefficients extracted from the same frozen window;
\item cross-refinement distance convergence using the successive Euclidean distance between adjacent coefficient vectors $(b_0,b_1,b_2)$;
\item a deterministic generic-graph control hierarchy with the same total node counts, realized as four repeated scalar open-grid Laplacian blocks with open boundary connectivity.
\end{enumerate}

No stochastic sampling is used anywhere in the phase.

\section{Frozen Branch Coefficients}

Phase IX tracks the same size-factored surrogate already frozen in Phase VIII:
\[
T_h(t) \approx h^{-2}\bigl(b_0(h) + b_1(h)t + b_2(h)t^2\bigr).
\]

\begin{table}[t]
\centering
\small
\begin{tabular}{@{}rrrrrr@{}}
\toprule
Level & $h$ & $b_0(h)$ & $b_1(h)$ & $b_2(h)$ & $R^2$ \\
\midrule
R1 & 0.083333 & 3.992541196696 & -15.196408636509 & 27.978374055231 & 0.999996635281 \\
R2 & 0.041667 & 3.992270826129 & -15.383158273830 & 28.598645942437 & 0.999996461490 \\
R3 & 0.027778 & 3.992219728715 & -15.417962259342 & 28.714851153144 & 0.999996428345 \\
R4 & 0.020833 & 3.992201767437 & -15.430159990250 & 28.755622542545 & 0.999996416672 \\
\bottomrule
\end{tabular}
\caption{Phase IX scaled coefficients on the frozen branch.}
\label{tab:coefficients}
\end{table}

The relative spans are:
\[
\mathrm{span}(b_0)=0.000085020802,\quad
\mathrm{span}(b_1)=0.015221236998,\quad
\mathrm{span}(b_2)=0.027260519706.
\]
The successive coefficient-vector distances are
\[
0.647775,\;0.121305,\;0.042557,
\]
which shrink monotonically across refinement. All three coefficients prefer the same simple extrapolation family, linear in $h^2$, with high fit quality.

\section{Ratio Stabilization}

The first normalized descriptors are the dimensionless ratios
\[
\frac{b_1(h)}{b_0(h)}, \qquad \frac{b_2(h)}{b_0(h)}.
\]

\begin{table}[t]
\centering
\small
\begin{tabular}{@{}lrrrr@{}}
\toprule
Descriptor & R1 & R2 & R3 & R4 \\
\midrule
$b_1/b_0$ & -3.806200 & -3.853235 & -3.862002 & -3.865075 \\
$b_2/b_0$ & 7.007661 & 7.163503 & 7.192703 & 7.202948 \\
\bottomrule
\end{tabular}
\caption{Dimensionless ratio descriptors on the frozen branch.}
\label{tab:ratios}
\end{table}

Their relative spans are
\[
0.015305769735 \quad \text{and} \quad 0.027344664673,
\]
both far smaller than the largest raw-coefficient span inherited from the unscaled trace fit, which is approximately $2.0077$. This is the first main Phase IX stabilization signal: the ratios are more stable than the raw coefficients.

The subset sensitivity remains bounded. Using the front four or rear four probe times of the frozen Phase VIII window changes the ratios by at most $0.06968$ relative to the full-window values. The deterministic truncation probes also remain bounded, with maximum ratio drift $0.07203$ for $\lambda \le 7.5$ and smaller drift for $\lambda \le 7.8$.

\section{Rescaled Invariant-Like Descriptors}

Phase IX next tracks the rescaled raw-coefficient descriptors
\[
I_k(h) = h^p a_k(h),
\]
with $p = 1.972182228379$ fixed from Phase VII.

\begin{table}[t]
\centering
\small
\begin{tabular}{@{}lrrrr@{}}
\toprule
Descriptor & R1 & R2 & R3 & R4 \\
\midrule
$I_0$ & 4.278286 & 4.361284 & 4.410698 & 4.446117 \\
$I_1$ & -16.284011 & -16.805054 & -17.034126 & -17.184576 \\
$I_2$ & 29.980777 & 31.242075 & 31.724840 & 32.025149 \\
\bottomrule
\end{tabular}
\caption{Rescaled invariant-like descriptors on the frozen branch.}
\label{tab:invariants}
\end{table}

Their relative spans are
\[
0.038369228803,\quad 0.053519229865,\quad 0.065434108792.
\]
These values are not flat in any literal sense, but they remain compact and refinement ordered. The largest subset-induced drift is $0.07057$, and the largest truncation-induced drift is $0.11088$. This is enough for a bounded invariant-tracking statement, not for a geometric interpretation.

\section{Generic-Graph Control Comparison}

Phase IX uses one deterministic control hierarchy only: a same-node-count open-grid scalar-block control built from four repeated open-boundary scalar Laplacian blocks. It keeps the same total dimensions as the frozen branch, but its connectivity pattern differs.

\begin{table}[t]
\centering
\small
\begin{tabular}{@{}lrrr@{}}
\toprule
Control coefficient & relative span & best extrapolation & limit estimate \\
\midrule
$b_0$ & 0.000302231342 & linear in $h$ & 3.992101755014 \\
$b_1$ & 0.078116509612 & linear in $h$ & -15.517159116109 \\
$b_2$ & 0.121073848083 & linear in $h$ & 29.012971564412 \\
\bottomrule
\end{tabular}
\caption{Control-hierarchy coefficient summary.}
\label{tab:control}
\end{table}

The separation is clear:
\begin{itemize}
\item control max ratio span: $0.121364602609$;
\item branch max ratio span: $0.027344664673$;
\item ratio-span multiplier (control / branch): $4.4383$;
\item final coefficient-vector distance multiplier (control / branch): $8.1097$;
\item branch best extrapolation family: linear in $h^2$ for all three scaled coefficients;
\item control best extrapolation family: linear in $h$ for all three scaled coefficients.
\end{itemize}

This is the second main Phase IX signal: the branch is not just stabilizing in isolation; it stabilizes differently from the deterministic generic control.

\section{Bounded Interpretation}

Phase IX supports the following narrow authority statement:
\begin{quote}
The frozen branch-local hierarchy supports reproducible coefficient stabilization, ratio stabilization, and invariant-like refinement tracking inside the already frozen Phase VIII short-time trace contract.
\end{quote}

This statement is justified by four combined facts:
\begin{enumerate}
\item at least one scaled coefficient, $b_0$, shows an immediate plateau-level span;
\item the dimensionless ratios are materially more stable than the raw trace coefficients;
\item the rescaled descriptors remain in compact refinement bands under the frozen probe subsets and truncation checks;
\item the deterministic open-grid control behaves differently in both span magnitude and extrapolation family.
\end{enumerate}

These facts are enough to freeze the Phase IX descriptor layer for later phases. They are not enough to claim a continuum invariant, a geometric coefficient, or a universal law.

\section{Non-Claims}

This paper does not claim:
\begin{enumerate}
\item continuum coefficients;
\item geometric interpretation of any descriptor;
\item universal invariants;
\item physical correspondence or law;
\item any continuum bridge.
\end{enumerate}

\section{Conclusion}

Phase IX closes with a bounded stabilization result on the frozen branch. The Phase VIII coefficient-like structures do not drift arbitrarily under refinement; their scaled versions compactify, their normalized ratios stabilize more strongly than the raw coefficients, their rescaled descriptors remain in bounded refinement bands, and the frozen branch is cleanly distinguishable from a deterministic generic-graph control. This is sufficient to freeze the Phase IX descriptor layer inside the spectral-trace contract for later stages.

Phase IX establishes coefficient stabilization and invariant-tracking feasibility for the frozen operator hierarchy.

\end{document}
